% =========================================================================================
% Chapter 5 -- Data Preparation (CRISP-DM)
%
% Split out of the former combined data chapter. The cut point is deliberate: this file
% begins where the description of the data ends and the operations on it begin. No
% paragraph was reordered in the split.
% =========================================================================================
\section{Data Preparation}
\label{sec:dataprep}

\subsection{Introduction}

Chapter~\ref{sec:data} established what these datasets contain and what counts as one
independent observation in each. This chapter describes what was then done to them: the
exclusions applied, the distributions that remained afterwards, the baselines those
distributions imply, and the partitioning that keeps a case on one side of the split.

The order matters. Every count reported here is a count after exclusion, and every baseline
is computed on the partition the module is actually evaluated on rather than on the pooled
data --- which is why they appear after the cleaning that produced them and not before.

\subsection{Cleaning and Exclusions}
\label{subsec:data-cleaning}

\paragraph{Lung.} Eleven files were excluded as off-target (not showing lung) and one as
flagged unusable by the dataset curators, reducing 198 files to the 187 used.

\paragraph{Gallbladder.} One class was found, on inspection, not to match its label. The
error is worth recording because of how it was made rather than merely that it occurred: the
class was first judged from a handful of images and appeared unremarkable. Viewing one image
per case, across all cases, showed otherwise.

\begin{keystatement}
Images from one case are not a sample of a class. They are a sample of a patient.
\end{keystatement}

The rule generalises beyond this incident and was applied to every subsequent inspection.
The affected class was excluded from training; the reasoning is given in
Section~\ref{subsec:gallbladder-module}.

\paragraph{Triage.} Vital signs recorded as zero are not physiological readings but an
encoding of ``not measured'', and were converted to genuine missing values. Fields knowable
only after triage --- final diagnosis, waiting time, disposition --- were removed, since a
model that sees them predicts the past rather than the future.

\subsection{Class Distribution and Baselines}
\label{subsec:data-distribution}

No score means anything on its own. A result must be read against what could be achieved
without a model at all, and the appropriate reference depends on the metric.

\begin{itemize}[label=$\bullet$]
  \setlength\itemsep{0.2em}
  \item For plain \textbf{accuracy}, the reference is the \textit{majority-class baseline}:
    the score obtained by always predicting the commonest class, ignoring the patient
    entirely.
  \item For \textbf{balanced accuracy} --- the mean of the per-class recalls --- that
    strategy scores no better than guessing, so the reference is the \textit{chance level},
    \(1/K\) for \(K\) classes.
  \item For \textbf{AUROC}, the reference is 0.500 regardless of how imbalanced the classes
    are, since the measure is invariant to prevalence.
\end{itemize}

Tab.~\ref{tab:baselines} gives the reference for each module.

\begin{table}[H]
  \centering
  \small
  \begin{tabular}{lllr}
    \toprule
    \textbf{Module} & \textbf{Metric reported} & \textbf{Reference} &
    \textbf{Value} \\
    \midrule
    Triage      & accuracy          & majority class, test split & 40.0\% \\
    Triage      & balanced accuracy & chance, 3 classes          & 33.3\% \\
    Gallbladder & balanced accuracy & chance, 5 classes          & 20.0\% \\
    Gallbladder & accuracy          & majority class             & 33.2\% \\
    Lung        & AUROC             & chance                     & 0.500 \\
    \bottomrule
  \end{tabular}
  \caption{The reference each module's results must be read against. Baselines are computed
    on the partition the module is evaluated on, not on the pooled data.}
  \label{tab:baselines}
\end{table}

One subtlety in the triage figures deserves stating, because using the wrong number would
misrepresent the result in both directions. Across the pooled data, 95{,}234 of 164{,}700
visits are \textit{high} urgency, giving a majority baseline of 57.8\%. But the model is
trained on earlier years and evaluated on the most recent, and in that test partition the
class mix is different: the majority baseline there is 40.0\%. The test-set figure is the
one a reported accuracy must be compared against.

The class imbalance is real in every module. At case level the gallbladder classes range
from 66 cases (inflammatory conditions) down to 25 (wall thickening) out of 199. In the lung
data the four findings appear in 63, 60, 46 and 24 of 187 clips --- so pleural effusion
rests on 24 positive examples, which is why it is later marked unreliable rather than
reported as a result.

One finding was dropped entirely. Pneumothorax was in the intended label set but has zero
positive examples in the usable data, so the lung module reports four findings rather than
five. A class with no positives cannot be learned, and a model claiming to detect it would
be claiming something it had never seen.

\subsection{Partitioning and Leakage Prevention}
\label{subsec:data-splits}

\paragraph{Imaging modules.} All three use grouped cross-validation by case: every file
belonging to one case is assigned to the same fold, so no fold is evaluated on a patient it
was trained on. Calibration data is held out separately, because fitting a calibrator on the
same predictions used to evaluate it would make the resulting confidence circular.

\paragraph{Triage.} The split is temporal rather than random --- earlier years for training,
the most recent year for testing --- giving 148{,}520 training and 16{,}180 test visits. A
random split would let the model learn from visits that occurred after the ones it is tested
on, which no deployed system could do.

\paragraph{A leak found by inspection, not by metrics.} Pooling four sources initially
produced a markedly higher accuracy than any single source, which is the wrong direction for
a harder, more heterogeneous problem. Investigation showed one source recorded vital signs
almost exclusively for a particular urgency grade, as an artefact of its collection process.
The model was not learning that abnormal vitals indicate urgency; it was learning that the
\emph{presence of a vitals record} indicates a label.

The fix was to null the affected fields for that source, keeping only those independently
reliable. The cost is visible in the data: roughly 87--91\% of the pooled visits now carry no
vital signs at all. That is a large loss of information, accepted because the alternative was
a higher number that meant nothing.

\subsection{Limitations}
\label{subsec:data-limits}

Stated as bounds on what may later be claimed:

\begin{itemize}[label=$\bullet$]
  \setlength\itemsep{0.2em}
  \item \textbf{Case counts are small.} 220 gallbladder cases across nine classes, and 165
    lung cases, are the true sample sizes. Confidence intervals on per-class results are
    correspondingly wide.
  \item \textbf{The lung case mapping is inferred.} Residual leakage is possible, and can
    only act in the optimistic direction.
  \item \textbf{The sources are public collections, not consecutive clinical series.} They
    are not a random sample of any population, so prevalence in these data implies nothing
    about prevalence in an emergency department.
  \item \textbf{Most vital-sign data was discarded.} The triage model operates largely on
    demographics and presenting complaint, which bounds how well it could possibly perform.
  \item \textbf{No external validation cohort} was available for the lung, gallbladder or
    triage modules. Only the cardiac module was tested on data from a different source.
\end{itemize}

\subsection{Conclusion}

The work described here produced no model and improved no metric. It did the opposite: by
establishing that the unit of analysis is the case rather than the file, it reduced the
apparent size of every imaging dataset --- most dramatically from 10{,}694 to 220 --- and
lowered every score that followed.

That is the point. The numbers reported in Chapter~\ref{sec:perception} are smaller than
those a file-level evaluation would have produced, and they are the ones that describe how
the models behave on a patient they have not seen.
